\subsection{Projective modules, injective modules, and extension modules}

PROPOSITION 10. --- Let M be an A-module. The following conditions are equivalent:

\begin{enumerate}
\def\labelenumi{(\roman{enumi})}
\item
  M is projective.
\item
  \(\operatorname{Ext}_{\mathbf{A}}^{i}(\mathbf{M},\mathbf{N}) = 0\) for every A-module \(\mathbf{N}\) and for every integer \(i > 0\) .
\item
  \(\mathrm{Ext}_{\mathbf{A}}^{1}(\mathbf{M},\mathbf{N}) = 0\) for every A-module \(\mathbf{N}\) .
\item
  There exists an exact sequence
\end{enumerate}

\[
0 \rightarrow \mathrm{K} \xrightarrow {u} \mathrm{P} \xrightarrow {r} \mathrm{M} \rightarrow 0,
\]

\(\text{where P is projective, and where } \text{ Ext}_{\mathbf{A}}^{1}(\mathbf{M}, \mathbf{K}) = 0.\)

\begin{enumerate}
\def\labelenumi{(\roman{enumi})}
\item
  \(\Rightarrow\) (ii): this is the corollary of prop. 5 of X, p. 88.
\item
  \(\Rightarrow\) (iii): this is trivial.
\item
  \(\Rightarrow\) (iv): this is clear since M is a quotient of a free module P.
\item
  \(\Rightarrow\) (i): since \(\operatorname{Ext}_{\mathbf{A}}^{1}(\mathbf{M},\mathbf{K}) = 0\) the canonical map
\end{enumerate}

\[
\operatorname{Hom} _ {\mathbf {A}} (\mathbf {M}, \mathbf {P}) \rightarrow \operatorname{Hom} _ {\mathbf {A}} (\mathbf {M}, \mathbf {M})
\]

is surjective (X, p. 90, corollary); there therefore exists an A-linear section of v and M is isomorphic to a direct factor of P, hence is projective.

PROPOSITION 11. --- Let N be an A-module. The following conditions are equivalent:

\begin{enumerate}
\def\labelenumi{(\roman{enumi})}
\item
  N is injective.
\item
  \(\operatorname{Ext}_{\mathbf{A}}^{i}(\mathbf{M},\mathbf{N}) = 0\) for every A-module M and every integer \(i > 0\) .
\item
  \(\operatorname{Ext}_{\mathbf{A}}^{1}(\mathbf{M},\mathbf{N}) = 0\) for every A-module \(\mathbf{M}\) .
\item
  There exists an exact sequence
\end{enumerate}

\[
0 \rightarrow \mathrm{N} \xrightarrow {u} \mathrm{I} \xrightarrow {v} \mathrm{C} \rightarrow 0,
\]

\(\text{where I is injective and where } \text{ Ext}_{\mathbf{A}}^{1} (\mathbf{C}, \mathbf{N}) = 0;\)

\begin{enumerate}
\def\labelenumi{(\alph{enumi})}
\setcounter{enumi}{21}
\item
  \(\operatorname{Ext}_{\mathbf{A}}^{1}(\mathbf{M},\mathbf{N}) = 0\) for every monogenic A-module \(\mathbf{M}\) .
\item
  \(\Rightarrow\) (ii): this is the corollary of prop. 5 of X, p. 88.
\end{enumerate}

\begin{enumerate}
\def\labelenumi{(\roman{enumi})}
\setcounter{enumi}{1}
\item
  \(\Rightarrow\) (iii) \(\Rightarrow\) (v): this is trivial.
\item
  \(\Rightarrow\) (iv): this is clear since N is a submodule of an injective module (X, p. 19, cor. 3).
\item
  \(\Rightarrow\) (i): since \(\operatorname{Ext}_{\mathbf{A}}^{1}(\mathbf{C},\mathbf{N}) = 0\) , the canonical homomorphism
\end{enumerate}

\[
\operatorname{Hom} _ {\mathbf {A}} (\mathrm{I}, \mathbf {N}) \rightarrow \operatorname{Hom} _ {\mathbf {A}} (\mathbf {N}, \mathbf {N})
\]

is surjective (X, p. 93, corollary); there therefore exists an A-linear retraction of u and N is isomorphic to a direct factor of I, hence is injective (X, p. 16, prop. 9).

\begin{enumerate}
\def\labelenumi{(\alph{enumi})}
\setcounter{enumi}{21}
\tightlist
\item
  \(\Rightarrow\) (i): if a is an ideal of A, one has \(\operatorname{Ext}_{\mathbf{A}}^{1}(\mathbf{A}/\mathfrak{a},\mathbf{N})=0\) ; the canonical map \(\operatorname{Hom}_{\mathbf{A}}(\mathbf{A},\mathbf{N})\to\operatorname{Hom}_{\mathbf{A}}(\mathfrak{a},\mathbf{N})\) is therefore surjective and N is injective (X, p. 16, prop. 10).
\end{enumerate}

